\begin{textsample}{POS Dim 3 – df – Score 21.00 – df/df\_inf\_4.txt}  \label{ex:f3_pos_123}
4.
EDUCAÇÃO \textbf{INFANTIL} PARA QUÊ?
A Educação \textbf{Infantil} é duplamente protegida pela Constituição Federal – CF ( BRASIL, 1988 ) : tanto é direito das \textbf{crianças} com idade entre zero e cinco anos ( Art. 208, IV ), como é direito das trabalhadoras e dos trabalhadores das cidades e do campo em relação às suas filhas, filhos e dependentes ( Art. 7, XXV ).
Ou seja, a Educação \textbf{Infantil} ilustra a relação recíproca que caracteriza os direitos humanos ao unir o direito à educação e ao trabalho.
Nesse sentido, a Educação \textbf{Infantil} \textbf{volta} -se como expressão dos direitos humanos, com foco na dignidade e no direito de aprendizagem das \textbf{crianças}.
Além disso, representa possibilidades de emancipação, uma vez que a garantia de oferta da Educação \textbf{Infantil} viabiliza o ingresso ou permanência de trabalhadoras e trabalhadores, com destaque às mulheres, no mercado de trabalho.
É importante mencionar que nem sempre o atendimento educativo foi assegurado às \textbf{crianças} dessa faixa \textbf{etária}.
Voltando um pouco na história, até o século XVIII, não existia nem mesmo um olhar às singularidades das \textbf{crianças}.
Entre os séculos XIX e XX, começam a emergir, de modo ainda incipiente, uma inflexão na direção dos direitos das \textbf{crianças}, prerrogativas de cidadania, teorias do desenvolvimento e conhecimento da periodicidade da vida \textbf{infantil}.
Iniciativas da Medicina, da Psicologia e da Pedagogia formulam discursos e sustentam práticas, divulgando normas de \textbf{higiene} e \textbf{cuidados} para as \textbf{crianças}.
Investem -se em campanhas de amamentação, criam -se \textbf{instituições} de atendimento, como as \textbf{creches} ( ainda com um viés assistencial ) e jardins de \textbf{infância}.
Com o tempo, começa a ganhar corpo um ideário sobre a \textbf{infância} que atribui à \textbf{criança} o status de sujeito de direitos, estendendo -se na elaboração de documentos celebrados internacionalmente, entre os quais, a Declaração de Genebra ( 1924 ), a Declaração Universal dos Direitos da Criança ( 1959 ) e a Convenção dos Direitos da Criança ( 1989 ).
No Brasil, a década de 1980 \textbf{marca} a virada do processo de reconhecimento e valorização da \textbf{infância}, porque o enfoque sai da tutela da família e recai sobre o direito assegurado pelo Estado.
A \textbf{criança} passa a ser considerada sujeito de direitos, fruto da mobilização da sociedade civil organizada, do movimento de mulheres e de pesquisadoras e pesquisadores da educação, em especial da Educação \textbf{Infantil}, que, por meio de intensas lutas e discussões sobre a necessidade da educação formal, culminou com os avanços registrados na CF de 1988, que passa a considerar a \textbf{criança} como sujeito de direitos : direito à vida, à saúde, à alimentação, à educação, ao lazer, à dignidade, ao respeito, à liberdade, às convivências familiar e comunitária.
Uma das consequências desse movimento é o reconhecimento da Educação \textbf{Infantil} como dever do Estado e direito da \textbf{criança}.
Se a promulgação do Estatuto da Criança e do Adolescente – ECA em 1990 foi um dos primeiros marcos nessa direção, a Lei de Diretrizes e Bases da Educação Nacional – LDB, promulgada em dezembro de 1996, é a consolidação que firma o elo entre a primeira \textbf{infância} e o atendimento educativo em \textbf{instituições} de educação coletiva.
A Educação \textbf{Infantil}, segundo os artigos 29 e 30 da LDB é a “ primeira etapa da Educação Básica ”.
Essa lei consagra definitivamente o atendimento às \textbf{crianças} de até cinco anos de idade, como parte da estrutura e do funcionamento dos sistemas educacionais.
Seguindo a mesma direção, a BNCC define o conjunto orgânico e progressivo de aprendizagens essenciais para a Educação \textbf{Infantil} e demais etapas da Educação Básica, afirmando a necessidade e importância de atendimento educativo às \textbf{crianças} da primeira \textbf{infância}. > A Base Nacional Comum Curricular ( BNCC ) é um documento de caráter normativo que define o conjunto orgânico e progressivo de aprendizagens essenciais que todos os alunos devem desenvolver ao longo das etapas e modalidades da Educação Básica, de modo a que tenham assegurados seus direitos de aprendizagem e desenvolvimento, em conformidade com o que preceitua o Plano Nacional de Educação ( BRASIL, 2017, p. 05 ).
A Educação \textbf{Infantil} organiza -se em dois momentos, denominados \textbf{Creche} e Pré-escola.
Tais denominações são controversas.
A história da Educação \textbf{Infantil} no Brasil tem se pautado numa luta entre superar o assistencialismo, por muito tempo associado à \textbf{creche}, e a preparação para o Ensino Fundamental, também, por algum tempo, ligada à \textbf{pré-escola}.
Dessa forma, quando se fala em \textbf{Creche} e Pré-escola, não se vincula a nenhuma dessas concepções ; trata -se, na verdade, da organização da primeira etapa da Educação Básica.
Uma nova organização dentro dessa já estabelecida na legislação brasileira foi apresentada pela BNCC : \textbf{bebês} ( de 0 a 1 ano e 6 \textbf{meses} ), \textbf{crianças} bem \textbf{pequenas} ( de 1 ano e 7 \textbf{meses} a 3 anos e 11 \textbf{meses} ) e \textbf{crianças} \textbf{pequenas} ( de 4 anos a 5 anos e 11 \textbf{meses} ), compreendendo esses três períodos singulares da \textbf{infância} em suas especificidades e necessidades para cada momento do desenvolvimento, sem a pretensão de enturmação seriada, que tem como critério as idades estanques.
Entende -se essa forma de organização como constituinte da unidade da Educação \textbf{Infantil} – Primeiro Ciclo, segundo a organização da Educação Básica da SEEDF.
A Educação \textbf{Infantil} não é assistencial, tampouco preparatória, pois trata -se de uma etapa da Educação Básica que abarca os direitos de aprendizagem voltados às reais e atuais necessidades e interesses das \textbf{crianças}, no sentido de proporcionar seu desenvolvimento integral.
Segundo o artigo 29 da LDB, a Educação \textbf{Infantil} tem como finalidade “ o desenvolvimento integral da \textbf{criança} até cinco anos em seus aspectos físico, psicológico, intelectual e social, complementando a ação da família e comunidade ”.
E, conforme o artigo 5º das Diretrizes Curriculares Nacionais para a Educação \textbf{Infantil} – DCNEI de 2010, a Educação \textbf{Infantil} é oferecida em estabelecimentos de educação, que se caracterizam como espaços institucionais não \textbf{domésticos}.
Esses estabelecimentos são públicos ou privados e precisam educar \textbf{cuidando} e \textbf{cuidar} educando, compreendendo a unidade indissociável desses Eixos Integradores, entre \textbf{crianças} de zero a cinco anos e onze \textbf{meses} de idade no período diurno, em jornada integral ou parcial.
Em seu artigo 8º, as DCNEI ressaltam que o objetivo principal da primeira etapa da Educação Básica é colaborar para o desenvolvimento integral das \textbf{crianças} ao garantir aprendizagens, bem como o direito à proteção, à saúde, à liberdade, ao respeito, à dignidade, à \textbf{brincadeira}, à convivência e à interação com \textbf{crianças} de diferentes faixas \textbf{etárias} e com \textbf{adultos}.
As perspectivas crítica e pós-crítica compreendidas nos pressupostos teóricos do Currículo em Movimento, como também a Psicologia Histórico-Cultural e Pedagogia Histórico-Crítica, apresentam o ato educativo como profundamente revolucionário, no sentido de provocar nas pessoas mudança de vida a partir da apropriação do patrimônio cultural da humanidade.
Nas interações, por meio do uso de instrumentos e signos, as pessoas se humanizam, são modificadas pela cultura e a modificam, numa relação dialética.
Tais perspectivas enfatizam também a constituição da individualidade a partir da coletividade.
Dessa forma, por meio das interações e \textbf{brincadeiras}, ocorre a vivência das práticas sociais, contempladas pelos campos de experiência e a apropriação dos saberes necessários, o que provocará uma nova formação.
É importante lembrar que Vigotski ( 2012a ) apresenta uma periodização das idades que não é estanque, pois depende das experiências culturais estabelecidas.
A cada nova idade ( ou período ), a \textbf{criança} vivencia experiências que contribuem para novas formações.
Estas inauguram e apontam transformações psicológicas, bem como geram uma nova situação social do desenvolvimento.
Essencialmente, essas teorias entendem que cada ser humano é diferente, portanto, segue caminhos diversos para aprender e desenvolver -se.
Assim, estruturar um currículo sobre essas bases implica lançar \textbf{mão} de práticas pedagógicas inovadoras e abertas, que proporcionem as descobertas, o respeito ao momento do desenvolvimento e às necessidades de cada ser humano e, no que diz respeito à primeira \textbf{infância}, que proponham ações educativas com intencionalidade a fim de fomentar o desenvolvimento da criatividade, da colaboração intra e intergeracional, da imaginação e da participação, enfatizando os princípios éticos, estéticos e políticos sobre os quais se fundamentam a Educação \textbf{Infantil} ( BRASIL, 2010a ).
A constituição da sociedade deve ser permeada pelo pleno respeito às \textbf{crianças}, em constante processo de valorização do protagonismo \textbf{infantil}, com a garantia de diferentes formas de sua participação, tanto no planejamento como na realização e avaliação das atividades que elas participam no contexto da \textbf{instituição} que oferta Educação \textbf{Infantil}.
No ano de 2013, foi instituída a Lei Federal nº 12.796/2013 que altera a LDB 9.394/1996 e determina que a educação obrigatória e gratuita atenda às \textbf{crianças} e adolescentes de quatro a 17 anos de idade, resultando na obrigatoriedade de as famílias e/ou responsáveis matricularem suas \textbf{crianças} na Educação \textbf{Infantil} a partir da idade estabelecida.
Na SEEDF, o atendimento educativo às \textbf{crianças} entre zero e cinco anos e onze \textbf{meses} de idade encontra -se distribuído em unidades públicas diversas : Jardim de Infância – JI e Centro de Educação \textbf{Infantil} – CEI.
Também, por questões estruturais, tem -se ainda turmas de Educação \textbf{Infantil} em espaços que não atendem às especificidades das \textbf{crianças}, como Escola Classe – EC, Centro de Atendimento Integral à Criança – CAIC, Centro de Ensino Fundamental – CEF, Centro Educacional – CED.
Além disso, a SEEDF estabelece parcerias com Organizações da Sociedade Civil – OSC, que são as \textbf{Instituições} Educacionais \textbf{Parceiras}.
Essas \textbf{instituições} atendem em \textbf{prédios} próprios ou públicos construídos pela SEEDF em parceria com o Ministério da Educação – MEC, que são as unidades do Pro-Infância, denominadas no Distrito Federal de Centros de Educação da Primeira Infância – CEPI.
Outro grupo a ser considerado é o das \textbf{instituições} privadas, confessionais ou não, que apresentam organizações diversas.
A maior oferta desta Secretaria concentra -se no atendimento educativo às \textbf{crianças} \textbf{pequenas} ( de 4 anos a 5 anos e 11 \textbf{meses} ).
Já em relação aos \textbf{bebês} ( de 0 a 1 ano e 6 \textbf{meses} ) e às \textbf{crianças} bem \textbf{pequenas} ( de 1 ano e 7 \textbf{meses} a 3 anos e 11 \textbf{meses} ), ainda há a necessidade de estender o atendimento educativo visando oportunizar a Educação \textbf{Infantil} para todas as faixas \textbf{etárias}, como prevê tanto a Meta 1 do Plano Nacional de Educação – PNE ( 2014-2024 ), como a Meta 1 do Plano Distrital de Educação – PDE ( 2015-2024 ).
Essa política de expansão deve vincular a garantia de qualidade na infraestrutura e equipamentos, na gestão, na formação dos profissionais da educação e nas práticas pedagógicas e avaliativas.
Em relação ao tempo de permanência das \textbf{crianças} nas \textbf{instituições} de educação coletiva, oferta -se tanto jornada de tempo parcial ( cinco horas ), quanto de tempo integral ( dez horas ).
Em todos os casos, ressalta -se que os profissionais devem trabalhar pela promoção das aprendizagens e do desenvolvimento integral das \textbf{crianças}.

% matched lemmas: adulto, bebê, brincadeira, creche, criança, cuidado, cuidar, doméstico, etário, higiene, infantil, infância, instituição, marca, mão, mês, parceiro, pequeno, pré-escola, prédio, volta
\end{textsample}
